\section{Анализ существующих подходов и форматов данных}

TODO2

TODO3

Качество решения в задачах машинного обучения на инженерной геометрии почти полностью определяется качеством и организацией набора данных. В проекте «НейроКОМПАС-3D» датасет имеет структуру последовательностей: каждой модели соответствует папка, содержащая шаги построения. Для каждого шага доступны файлы STEP и STL, а также файл текстового описания операций между шагами; в более поздней редакции пайплайна используются JSON-метаданные. Такой формат воспроизводит реальную логику CAD-построения, где объект существует как история преобразований, а не как статичная сетка.

Требования к формату данных были зафиксированы уже на первом этапе. Для каждой модели должны быть доступны идентификатор набора, номера шагов, файл с описанием операций и согласованный словарь поддерживаемых операций. В финальной реализации в типовой набор операций вошли: выдавливание (bossExtrusion), вырезание выдавливанием (cutExtrusion), скругление (fillet), фаска (chamfer). Эта ограниченная номенклатура была выбрана осознанно: она позволяла локализовать исследование на наиболее типичных и одновременно существенно различающихся по геометрическим признакам операциях, не пытаясь сразу охватить весь спектр CAD-команд.

Важным элементом начальной постановки было понимание масштабов данных. Уже в ранних отчетных материалах указывалась ориентировочная потребность порядка 10 000 моделей по 20–30 шагов для этапов 2 и 3. Практика полностью подтвердила эту оценку. Первичная тестовая выборка из 362 наборов оказалась достаточной для демонстрации обучаемости, но недостаточной для достоверного сравнения архитектур и устойчивых выводов по качеству. Следовательно, в проекте изначально корректно оценивалась высокая data hunger задачи.

Особое внимание в проекте уделялось препроцессингу и численной нормализации. На этапе 3 STL-файлы преобразовывались в point cloud размерности N×3, геометрия переносилась в окрестность начала координат и масштабировалась к единичному размеру. Steps.txt разбирался на пары «ключ-значение»; ключи кодировались one-hot-векторами; значения нормировались в зависимости от типа — углы переводились в 0–1, длины домножались на коэффициент масштаба, координаты сдвигались на дельту переноса.

Отдельной проблемой стала координатная система операции и ее связь с глобальной системой модели. В рабочих материалах прямо отмечалась необходимость хранить нормаль к плоскости эскиза либо матрицу поворота и смещение, позволяющие сопоставить систему STL и систему эскиза. Это замечание отражает важную теорему инженерного машинного обучения: если модель должна предсказывать параметры в локальной системе объекта, то либо вход должен быть инвариантен к глобальному положению, либо ему необходимо явно сообщить преобразование между системами координат. Игнорирование этого принципа ведет к резкому падению качества регрессии.

Применительно к сегментационному варианту датасет претерпел существенную переработку. STEP стал использоваться не столько как прямой вход для сети, сколько как источник точной геометрической семантики при формировании меток. Для пары соседних шагов определялись изменившиеся грани; далее для треугольников финальной модели вычислялось соответствие типу операции, породившей содержащую их грань. Тем самым была реализована масштабируемая схема автоматической разметки без ручного аннотирования каждого треугольника человеком.

Вторая нейронная линия относится к сегментации. В качестве базового варианта была рассмотрена MeshCNN-модель, работающая по ребрам и локальным отношениям смежности. Для сложных поверхностей такая архитектура привлекательна тем, что использует топологию сетки, а не только набор несвязанных точек. В перспективе также рассматривались PointCNN и иные модели, способные обобщать представление до point cloud и лучше переноситься на геометрию с большим числом треугольников. Фактически в проекте был исследован спектр между сеточно-зависимыми и сеточно-инвариантными представлениями.

